% ============================================================
%  MODULE V — Technical Analysis
% ============================================================

\chapter{Technical Analysis}

\begin{notebox}
Technical analysis (TA) is the study of \emph{price and volume history} to forecast future
price direction. It operates on the premise that all known information is reflected in the price,
and that crowd psychology creates recurring, identifiable patterns.
\end{notebox}

% ────────────────────────────────────────────────────────────
\section{Core Assumptions of Technical Analysis}

\begin{enumerate}[label=\textbf{\arabic*.}]
  \item \textbf{Price discounts everything} — market price reflects all fundamentals, sentiment,
        and information. Analysis focuses on price itself.
  \item \textbf{Prices move in trends} — once established, a trend tends to continue until it
        reverses. The trend is your friend.
  \item \textbf{History repeats} — patterns recur because human psychology is consistent.
        Fear and greed create the same structures again and again.
\end{enumerate}

% ────────────────────────────────────────────────────────────
\section{Chart Types}

\subsection{Candlestick Charts}
The most common chart type. Each candle represents a time period (1min, 1hr, 1day, etc.)

\begin{center}
\begin{tabular}{>{\bfseries}p{3cm} p{9.5cm}}
\toprule
Candle Part & Meaning \\
\midrule
Body & Range between Open and Close price \\
Upper Wick & Highest price reached during the period \\
Lower Wick & Lowest price reached during the period \\
Green / White body & Close $>$ Open (bullish period) \\
Red / Black body & Close $<$ Open (bearish period) \\
\bottomrule
\end{tabular}
\end{center}

\subsection{Key Candlestick Patterns}

\begin{center}
\begin{tabular}{>{\bfseries}p{3.5cm} p{3cm} p{6cm}}
\toprule
Pattern & Signal & Description \\
\midrule
Doji & Indecision & Open $\approx$ Close; wicks on both sides \\
Hammer & Bullish reversal & Small body at top, long lower wick \\
Shooting Star & Bearish reversal & Small body at bottom, long upper wick \\
Engulfing & Reversal & Second candle fully engulfs first \\
Marubozu & Strong trend & Full body, no wicks \\
\bottomrule
\end{tabular}
\end{center}

% ────────────────────────────────────────────────────────────
\section{Trend Analysis}

\subsection{Identifying Trends}
\begin{itemize}
  \item \textbf{Uptrend} — series of Higher Highs (HH) and Higher Lows (HL)
  \item \textbf{Downtrend} — series of Lower Highs (LH) and Lower Lows (LL)
  \item \textbf{Sideways / Ranging} — price oscillates between a ceiling and a floor
\end{itemize}

\subsection{Trendlines}
A line connecting a series of swing lows in an uptrend (or swing highs in a downtrend).
A break of a trendline is a potential signal of trend reversal.

\subsection{Support and Resistance}
\begin{keyterm}{Support}
A price level where buying interest is strong enough to prevent further decline. Price has
bounced from this level before. It is psychological ``floor''.
\end{keyterm}

\begin{keyterm}{Resistance}
A price level where selling pressure prevents further advance. Price has been rejected here
before. Acts as a psychological ``ceiling''.
\end{keyterm}

\begin{notebox}
\textbf{Role Reversal}: Once broken, a support level often becomes new resistance,
and a resistance level often becomes new support. This is one of the most reliable
observations in all of technical analysis.
\end{notebox}

% ────────────────────────────────────────────────────────────
\section{Moving Averages}

A moving average smooths price data to reveal the underlying trend direction.

\begin{formulabox}
\textbf{Simple Moving Average (SMA):}
\[
  SMA_n = \frac{P_1 + P_2 + \cdots + P_n}{n}
\]

\textbf{Exponential Moving Average (EMA):}
\[
  EMA_t = P_t \cdot k + EMA_{t-1} \cdot (1 - k), \quad k = \frac{2}{n+1}
\]
EMA gives more weight to recent prices; reacts faster than SMA.
\end{formulabox}

\subsection{Common MA Periods}
\begin{itemize}
  \item \textbf{20 / 21 EMA} — short-term trend; used by swing traders
  \item \textbf{50 SMA} — medium-term trend; key institutional reference
  \item \textbf{200 SMA} — long-term trend; major bull/bear dividing line
\end{itemize}

\subsection{Moving Average Crossovers}
\begin{itemize}
  \item \textbf{Golden Cross} — 50 SMA crosses above 200 SMA. Bullish signal.
  \item \textbf{Death Cross} — 50 SMA crosses below 200 SMA. Bearish signal.
\end{itemize}

% ────────────────────────────────────────────────────────────
\section{Key Technical Indicators}

\subsection{Relative Strength Index (RSI)}
Measures the speed and magnitude of price changes. Oscillates between 0 and 100.

\begin{formulabox}
\[
  RSI = 100 - \frac{100}{1 + RS}, \quad RS = \frac{\text{Average Gain (14 periods)}}{\text{Average Loss (14 periods)}}
\]
\end{formulabox}

\begin{itemize}
  \item RSI $>$ 70: Overbought (potential sell signal)
  \item RSI $<$ 30: Oversold (potential buy signal)
  \item Divergence: Price makes new high but RSI does not — bearish divergence.
\end{itemize}

\subsection{MACD (Moving Average Convergence Divergence)}
Two EMAs (12 and 26 periods) and a signal line (9-period EMA of MACD).
\begin{itemize}
  \item \textbf{MACD crosses above signal line} — bullish
  \item \textbf{MACD crosses below signal line} — bearish
  \item \textbf{Histogram} — visualises the gap between MACD and signal line
\end{itemize}

\subsection{Bollinger Bands}
A 20-period SMA with bands $\pm 2$ standard deviations above and below.
\begin{itemize}
  \item Price touching upper band: potentially overbought or strong trend continuation
  \item Price touching lower band: potentially oversold
  \item \textbf{Squeeze}: narrow bands signal low volatility — often precedes a big move
\end{itemize}

\subsection{Volume}
Volume validates price moves. High volume on a breakout = more convincing.
Low volume on a rally = suspect. Watch for \textbf{volume divergence}.

% ────────────────────────────────────────────────────────────
\section{Chart Patterns}

\subsection{Continuation Patterns}
\begin{itemize}
  \item \textbf{Flag / Pennant} — tight consolidation after a sharp move; usually resolves in trend direction
  \item \textbf{Triangle} (ascending, descending, symmetrical) — price coiling; breakout expected
  \item \textbf{Cup and Handle} — U-shaped recovery, small pullback, then breakout
\end{itemize}

\subsection{Reversal Patterns}
\begin{itemize}
  \item \textbf{Head and Shoulders} — three peaks; middle (head) is highest; neckline break = sell signal
  \item \textbf{Inverse Head and Shoulders} — same structure inverted; neckline break = buy signal
  \item \textbf{Double Top / Double Bottom} — price fails twice at same level; signals reversal
\end{itemize}

% ────────────────────────────────────────────────────────────
\section{Fibonacci Retracements}

Based on the Fibonacci sequence, key retracement levels are:
\textbf{23.6\%, 38.2\%, 50\%, 61.8\%, 78.6\%}

The 61.8\% level is the ``Golden Ratio'' and is considered the most significant.
Traders use these levels as potential support/resistance zones after a move.

\begin{warningbox}
Technical analysis is descriptive, not prescriptive. No indicator works all the time.
Use TA as a \emph{framework for probability}, not as a crystal ball. Always combine
with risk management; patterns fail regularly.
\end{warningbox}
